\section{Common domains and parameter-dependent coordinates}
\label{rw:app:domains}

Coefficient division compares a local map with its center value while
holding a common input fixed. A coefficient-dependent input change
therefore enters the coefficient derivative. We record this rule,
then check that phase changes preserve the whole passage and regular
connections preserve flat outputs. The physical domain is the
intersection of the scalar whole-path conditions, regular connection
domains, positive central heights and the chosen collar. The integral
comparisons use larger signed domains before this physical restriction.

\begin{lemma}[A fixed input under coefficient division]
\label{rw:app:fixed-input-division}
Let $b=(b_0,\eta)\in B_0\times\R^k$, and let a parameterized
input change $x=\phi(s,b)$ be defined on a common product $J\times B$.
Assume its inverse exists on a slightly larger common interval and both
maps have bounded derivatives through total order $N+1$.
Suppose $F(x,b)$ has those bounds on the entire image and
$F(x,b_0,0)=0$ for every represented center input. Define
$G(s,b)=F(\phi(s,b),b)$ and suppose
$(s,b_0,t\eta)$ remains in the common domain for $0\le t\le1$.
Then the actual divided columns are
\[
 G_i(s,b)=\int_0^1\partial_{\eta_i}G(s,b_0,t\eta)\,\dd t,
 \qquad G=\sum_i\eta_iG_i.
\]
They have bounded total order $N$ on $J\times B$. In general these
columns are not obtained by composing previously divided columns of $F$
with $\phi(s,b)$.
\end{lemma}
\begin{proof}
Apply the fundamental theorem of calculus to
$t\mapsto G(s,b_0,t\eta)$ with $s$ fixed. Its value at zero
vanishes by the center identity on the entire represented interval.
At each coefficient point the chain rule reads
\[
 \partial_{\eta_i}G
 =F_{\eta_i}(\phi,b)+F_x(\phi,b)\phi_{\eta_i}(s,b).
\]
Even though $F_x=0$ at the center, the second term need not vanish
along the coefficient segment. It must be retained in the divided
column.
Repeated chain and product rules through order $N$ use only the
bounded $N+1$ jets of $F$ and $\phi$ on the stated image. Integration
proves the bounds, and the prescribed inverse margins make this a
common coordinate on every open overlap.
\end{proof}

In Section~\ref{rw:joint-section}, form the separated columns only
after fixing the actual common input. Two input derivatives of these
columns use three input and one coefficient derivatives of the
nonconstant critical logarithm. Their coefficient Lipschitz bounds use
two of each. Both requirements have total order four and follow from
the order-five primitive estimate. Counting only derivatives at a fixed
coefficient would miss the second chain-rule term.

\begin{lemma}[Changing a scalar phase]
\label{rw:app:scalar-phase}
Suppose $x'=qg_1(x,b)$ and $y=\Phi(x,b)$ is an increasing regular
phase change, with positive bounded derivative and inverse derivative,
on a common entry-to-exit interval. Allow also a time change
$d\widetilde t/dt=a(x,b)>0$, with bounded positive unit $a$.
The exact new scalar factor is
\[
 g_2(\Phi(x,b),b)=\frac{\Phi_x(x,b)}{a(x,b)}g_1(x,b).
\]
The whole-path passing conditions are equivalent in the two phases.
For a normally contracting equation $Z'=-Z\beta(x,Z,b)$, the
linear normal action and the nonlinear relative integral are unchanged
when their time and phase factors are transformed together.
\end{lemma}
\begin{proof}
Differentiate $y=\Phi(x,b)$. Its positive multiplier and increasing
$\Phi$ preserve positivity on the whole closed entry-to-exit interval;
corresponding scalar zeros cannot be crossed in finite time. The normal
rate becomes $\widetilde\beta=\beta/a$, so change of variables gives
\[
 \int_{x_e}^{x_R}\frac{\beta(x,0,b)}{qg_1(x,b)}\,\dd x
 =\int_{\Phi(x_e)}^{\Phi(x_R)}
       \frac{\widetilde\beta(y,0,b)}{qg_2(y,b)}\,\dd y.
\]
The same substitution applies to the relative integral along the
actual orbit. At a separated lost exit, the normal domination and
integrable-tail bound of Appendix~\ref{rw:app:relative} justify taking
the infinite-time limit.
\end{proof}

The positive phase factors in Section~5 identify the entire passages
and their contractions.
For a quadratic factor a nonpositive gap retains only entries above
the positive root, toward the fixed right exit. For a weak factor
$u(\gamma+u)$ it retains $u>\max(0,-\gamma)$. For the cubic raw
factor $xF_I(x)$ one uses the positive component of $F_I$ adjoining
the fixed exit. Each condition must hold on the whole entry-to-exit
path; two endpoint values do not determine passage.

\begin{lemma}[Regular maps and flat outputs]
\label{rw:app:flat-composition}
Suppose a normally contracting transition has the form
$G(r,b)=L(b)r e^{j(r,b)}$, with bounded selected mixed derivatives
of $j$ and of the absolute factor $L$, while derivatives of
$\log L$ may be unbounded. Let two regular maps fixing zero have
positive bounded units
$\Phi_{\rm in}(r,b)=r a(r,b)$ and
$\Phi_{\rm out}(z,b)=z c(z,b)$. Their composition has the exact form
\[
 \Phi_{\rm out}\circ G\circ\Phi_{\rm in}
       =L a(0,b)c(0,b)r e^{\psi(r,b)},
\]
where the selected jets of $\psi$ are bounded on the common domain.
They extend to a lost passage whenever the relative unit and absolute
factor have the corresponding extensions. This statement supplies no
uniform inverse derivative for the composite map.
\end{lemma}
\begin{proof}
The exact relative logarithm is
\[
 \psi=\log\frac{a(r,b)}{a(0,b)}+j(ra(r,b),b)
 +\log\frac{c(Lra(r,b)e^{j(ra(r,b),b)},b)}{c(0,b)}.
\]
Every denominator has a fixed positive margin. The mixed derivatives
use absolute derivatives of $L$, so no factor $L^{-1}$ or coefficient
derivative of $\log L$ occurs. The chain rule gives the bounds and
extensions. At $L=0$ the composite and its input derivative vanish,
so the assertion gives no uniform inverse bound.
\end{proof}

For the contact and weak-center arguments, a displacement such as
$\delta\mapsto\delta+3(|\tau|+|\zeta|+|\nu|)$ chooses a center
with the original coefficient vector held fixed. Apply Taylor's formula
on that prescribed segment in the larger domain; the absolute values
are not differentiated. The center identity holds on an open passing
neighborhood of each retained input, so its passive center-gap
derivative vanishes. This gives the linear transverse term and the
quadratic remainder used in those arguments.
